% ============================================================
% BEISPIEL: Kapitel 1 - Einleitung (Stundenbuchungstool)
% Dies ist ein Beispiel, wie eine Einleitung aussehen koennte.
% ============================================================

\chapter{Problemstellung, Ziel und Vorgehensweise}
\label{chap:einleitung}

\section{Motivation und Ausgangssituation}
\label{sec:motivation}

In vielen Unternehmen stellt die korrekte Erfassung und Zuordnung von 
Arbeitsstunden auf Projekte eine grundlegende Voraussetzung für die 
Kosten\-rechnung und Projekt\-steuerung dar. Insbesondere in 
Forschungs- und Entwicklungsabteilungen, in denen Mitarbeitende 
parallel an mehreren Projekten arbeiten, ist eine präzise 
Stundendokumentation unerlässlich \cite[vgl.][S.~34]{controlling_buch}.

Die Digitalisierung solcher Prozesse bietet die Möglichkeit, manuelle 
Tätigkeiten zu reduzieren, die Datenqualität zu erhöhen und 
Auswertungen in Echtzeit bereitzustellen 
\cite[vgl.][S.~12\,ff.]{digitalisierung2024}. Gleichzeitig stellen 
Low-Code-Plattformen wie Microsoft Power Apps eine Technologie dar, 
die es ermöglicht, solche Anwendungen mit vergleichsweise geringem 
Entwicklungsaufwand zu realisieren \cite[vgl.][S.~8]{microsoft_lowcode}.

\section{Problemstellung}
\label{sec:problemstellung}

Der aktuelle Stundenbuchungsprozess beim Dualen Partner basiert auf 
einer Vielzahl individueller Excel-Dateien. Jede/r Mitarbeitende führt 
eine eigene Tabelle, in der die auf verschiedene Projekte verteilten 
Stunden dokumentiert werden. Zusätzlich existiert eine zentrale 
Projektliste mit Metadaten wie Projekt-ID, Projektname und 
Kostenstelle.

Dieser Prozess weist folgende Schwachstellen auf:

\begin{itemize}
    \item \textbf{Hoher manueller Aufwand:} Am Monatsende müssen 
    sämtliche Einzel-Tabellen manuell zusammengeführt werden.
    \item \textbf{Niedrige Datenqualität:} Durch Freitexteingaben 
    entstehen Tippfehler, inkonsistente Bezeichnungen und fehlende 
    Pflichtangaben.
    \item \textbf{Fehleranfälliger SAP-Import:} Die zusammengeführte 
    Tabelle muss in das Buchungssystem SAP importiert werden. 
    Auftretende Fehlermeldungen erfordern iterative manuelle 
    Korrekturen.
    \item \textbf{Eingeschränkte Auswertbarkeit:} Aufgrund der 
    heterogenen Datenstruktur sind Analysen wie 
    Stundensatzberechnungen nur mit erheblichem Aufwand möglich.
\end{itemize}

\section{Ziel der Arbeit}
\label{sec:ziel}

Ziel dieser Arbeit ist die Fertigstellung und Einführung eines 
digitalen Stundenbuchungstools, das den beschriebenen manuellen 
Prozess ablöst. Die Entwicklung erfolgt in Zusammenarbeit mit einer 
weiteren Entwicklerin, die bereits Vorarbeiten geleistet hat.

Der Eigenanteil dieser Arbeit gliedert sich in zwei Schwerpunkte:

\begin{enumerate}
    \item \textbf{Technische Umsetzung:} Analyse des bestehenden 
    Prozesses, systematische Ableitung der Anforderungen und 
    Implementierung der Kernfunktionen des Tools auf Basis von 
    Microsoft Power Apps und Power Automate.
    \item \textbf{Test- und Einführungsphase:} Erstellung eines 
    Testkonzepts mit definierten Testfällen, Überprüfung der 
    Datenqualität sowie Vorbereitung des Rollouts (Benutzerrollen, 
    Berechtigungen, Nutzerleitfäden).
\end{enumerate}

Nicht Gegenstand dieser Arbeit sind die Anbindung an weitere 
ERP-Module sowie die langfristige Wartung des Systems nach der 
Einführung.

\section{Geplantes Vorgehen}
\label{sec:vorgehen}

Das Vorgehen orientiert sich an einem iterativen Ansatz und lässt 
sich in folgende Phasen unterteilen:

\begin{enumerate}
    \item \textbf{Analyse des Ist-Zustands:} Untersuchung des 
    bestehenden Excel-basierten Prozesses hinsichtlich Schwachstellen, 
    Schnittstellen und impliziter Anforderungen.
    \item \textbf{Anforderungsdefinition:} Systematische Erfassung 
    funktionaler und nicht-funktionaler Anforderungen auf Basis der 
    Analyse und in Abstimmung mit den Stakeholdern.
    \item \textbf{Implementierung:} Technische Umsetzung der 
    Kernfunktionen in Microsoft Power Apps mit Automatisierungen 
    über Power Automate.
    \item \textbf{Test:} Erstellung und Durchführung eines 
    Testkonzepts zur Sicherstellung der Funktionsfähigkeit und 
    Datenqualität.
    \item \textbf{Einführungsvorbereitung:} Festlegung von 
    Benutzerrollen, Erstellung von Dokumentationen und Planung des 
    Rollouts.
\end{enumerate}

\section{Aufbau der Arbeit}
\label{sec:aufbau}

Die vorliegende Arbeit ist wie folgt aufgebaut: In 
Kapitel~\ref{chap:grundlagen} werden die theoretischen Grundlagen 
zu Low-Code-Entwicklung, Microsoft Power Platform und 
Datenqualitätsmanagement erläutert. 
Kapitel~\ref{chap:analyse} beschreibt die Analyse des Ist-Zustands 
und die daraus abgeleiteten Anforderungen. Die technische Umsetzung 
wird in Kapitel~\ref{chap:umsetzung} dokumentiert, gefolgt von der 
Test- und Einführungsphase in Kapitel~\ref{chap:test}. 
Abschließend erfolgt in Kapitel~\ref{chap:reflexion} eine kritische 
Reflexion der Ergebnisse sowie ein Ausblick auf weiterführende 
Maßnahmen.
